\documentclass[a4paper]{article}

% Español (México) con XeLaTeX: fontspec para fuentes Unicode, polyglossia
% para la división silábica y los nombres del documento en la variante mexicana.
\usepackage{fontspec}
\usepackage{polyglossia}
\setdefaultlanguage[variant=mexican]{spanish}
% Los nombres chinos van como «pinyin (original)»: xeCJK da a los caracteres
% Han su propia fuente; sin ella XeLaTeX los omite con un aviso.
% FandolSong no cubre algunos caracteres de nombres (U+73FA); WenQuanYi Zen Hei sí.
\usepackage[AutoFallBack=true]{xeCJK}
\setCJKmainfont{FandolSong-Regular.otf}
\setCJKfallbackfamilyfont{\CJKrmdefault}{WenQuanYi Zen Hei}
% polyglossia no traduce los nombres de listings.
\AtBeginDocument{\providecommand{\lstlistingname}{}\renewcommand{\lstlistingname}{Listado}%
  \providecommand{\lstlistlistingname}{}\renewcommand{\lstlistlistingname}{Índice de listados}}
\usepackage{amsmath,amssymb}
\usepackage{graphicx}
\usepackage[margin=2.25cm]{geometry}
\usepackage[most]{tcolorbox}
\usepackage{listings}
\usepackage{hyperref}
\usepackage{booktabs}
\usepackage{array}
\usepackage{float}
\usepackage{xcolor}
\usepackage{longtable}

\definecolor{kbBack}{HTML}{F2F6FA}\definecolor{kbFrame}{HTML}{9DB4CC}\definecolor{kbTitle}{HTML}{52708F}
\definecolor{ibBack}{HTML}{FBF5E7}\definecolor{ibFrame}{HTML}{C9A961}\definecolor{ibTitle}{HTML}{7D5F1E}
\definecolor{wbBack}{HTML}{FCF2F0}\definecolor{wbFrame}{HTML}{D6A096}\definecolor{wbTitle}{HTML}{9E4F43}

\tcbset{softbox/.style={enhanced,breakable,boxrule=0.8pt,arc=2pt,left=3mm,right=3mm,top=2mm,bottom=2mm,coltitle=white,fonttitle=\bfseries,toptitle=0.6mm,bottomtitle=0.6mm}}
\newtcolorbox{knowledgebox}[1]{softbox,colback=kbBack,colframe=kbFrame,colbacktitle=kbTitle,title={#1}}
\newtcolorbox{importantbox}[1]{softbox,colback=ibBack,colframe=ibFrame,colbacktitle=ibTitle,title={#1}}
\newtcolorbox{warningbox}[1]{softbox,colback=wbBack,colframe=wbFrame,colbacktitle=wbTitle,title={#1}}

\hypersetup{colorlinks=true,linkcolor=blue!50!black,urlcolor=blue!60!black}
\setlength{\parindent}{2em}
\setlength{\parskip}{0.35em}
\renewcommand{\arraystretch}{1.25}
\graphicspath{{./}}

\newcommand{\courseurl}{https://www.bilibili.com/video/BV1P4LX68EcS}
\newcommand{\sourcefigure}[4]{%
\begin{figure}[H]
\centering
\includegraphics[width=#1\textwidth]{#2}
\caption{#3\protect\footnotemark}
\end{figure}
\footnotetext{Intervalo del video: #4.}}

\begin{document}

\begin{titlepage}
\centering
\vspace*{0.6cm}
{\huge\bfseries Self-Evolving Agents 2026\par}
\vspace{0.3cm}
{\LARGE De la frontera académica a la implementación industrial\par}
\vspace{0.5cm}
{\large Apuntes de curso en chino del seminario académico temático de NICE\par}
\vspace{0.8cm}
\includegraphics[width=0.88\textwidth,height=0.44\textheight,keepaspectratio]{cover.jpg}
\vfill
\begin{tcolorbox}[width=0.93\textwidth,colback=black!2!white,colframe=black!55,sharp corners]
\textbf{Fuente}: NICE Academia, «Invitados globales conversan sobre Self-Evolving Agents: de la frontera académica a la implementación industrial»\par
\textbf{Fecha de publicación del video}: 2026-06-17\par
\textbf{Alcance del curso}: siete presentaciones temáticas, dos paneles; dos aperturas incorporadas en la introducción del curso\par
\textbf{Duración total del video original}: 5:54:29\par
\textbf{Enlace del video}: \href{\courseurl}{\nolinkurl{\courseurl}}\par
\textbf{Fecha de elaboración de las notas}: 2026-08-17\par
\textbf{Método de elaboración}: compilación integral source-first de las capturas del video público, la transcripción de Whisper y los artículos y materiales de proyecto publicados por los ponentes
\end{tcolorbox}
\end{titlepage}

\tableofcontents
\newpage

\section{Introducción al curso: qué evoluciona realmente un Agente que se automejora}

Este conjunto de presentaciones no limita la «automejora» a seguir ajustando los pesos del modelo, sino que concibe al Agente como un sistema compuesto en conjunto por \textbf{el modelo, el entorno, la experiencia, la memoria, las herramientas, los verificadores, la infraestructura de entrenamiento y los mecanismos de gobernanza}. Basta con que el sistema pueda sedimentar en decisiones posteriores la nueva información generada por la interacción para que ya cuente con cierta capacidad de mejora a lo largo del tiempo; la diferencia radica solo en si lo que se actualiza son los parámetros, el prompt, la memoria, las habilidades, el currículo del entorno, o todo el Agent harness.

\begin{importantbox}{El ciclo unificado que atraviesa las nueve sesiones}
Los métodos del curso pueden escribirse de manera unificada como: \textbf{plantear una tarea o un problema $\rightarrow$ interactuar con el entorno $\rightarrow$ obtener retroalimentación verificable $\rightarrow$ diagnosticar el modo de falla $\rightarrow$ actualizar la política, la experiencia o el entorno $\rightarrow$ entrar a la siguiente ronda de interacción}. Las divergencias entre las sesiones ocurren sobre todo en cuatro puntos: «de dónde viene el entorno», «si la retroalimentación es confiable», «qué estado se actualiza» y «cómo evitar la acumulación de experiencia errónea».
\end{importantbox}

\subsection{Cuatro dimensiones de análisis}

Al leer cada sesión conviene rastrear cuatro preguntas. Primera, \textbf{evolution target}: qué se actualiza en realidad, si los parámetros del modelo, el contexto, la memoria de largo plazo, la cadena de herramientas, la distribución de tareas o el entorno mismo. Segunda, \textbf{feedback source}: de dónde proviene la retroalimentación, si de verificadores basados en reglas, preferencias humanas, resultados de ejecución, el modelo del mundo, un Agente adversario o el mercado y la conducta de los usuarios. Tercera, \textbf{credit assignment}: a qué paso debe atribuirse el éxito o el fracaso dentro de una trayectoria larga. Cuarta, \textbf{safety boundary}: cómo evita el sistema que un éxito casual, una vulnerabilidad de la recompensa y una autoevaluación errónea se fijen como una nueva política.

\begin{knowledgebox}{Por qué organizar este conjunto de seminarios como un curso}
Las nueve unidades abordan el tema respectivamente desde el causal world model, el Agentic RL, la previsión de fallas, la síntesis de entornos, el aprendizaje a partir de la experiencia, la externalización, la teoría unificada y la práctica industrial. Ver una sola sesión de manera aislada hace fácil malinterpretar la automejora como el nombre de un algoritmo; solo al leerlas en el orden del curso se aprecia que en realidad se trata de un problema de diseño de sistemas que atraviesa la teoría del aprendizaje, la ingeniería de datos, la infraestructura de sistemas y el ciclo cerrado del producto.
\end{knowledgebox}

\subsection{Mapa de ruta del curso}

\begin{longtable}{p{0.08\textwidth}p{0.27\textwidth}p{0.26\textwidth}p{0.29\textwidth}}
\toprule
Sesión & Pregunta central & Mecanismo principal & Relación con el eje del curso \\
\midrule
1 & Cómo entiende el Agente un mundo abierto & causal world model, intervención activa, meta-causal graph & Aporta un modelo de entorno interpretable para la automejora \\
2 & Cómo producir Agent experience a escala & rollout infrastructure, entornos paralelos, reutilización recursiva de la experiencia & Convierte la experiencia de interacción en un recurso entrenable \\
3 & Por qué falla el foresight & desajuste del world model, entrenamiento explícito, curriculum generation & Pasa del diagnóstico de fallas a una mejora dirigida \\
4 & Cómo coevolucionan el Agente y el entorno & descubrimiento de entornos, síntesis de tareas, recompensa verificable, Agentic RL & Expande a la vez el entorno de entrenamiento y la capacidad de la política \\
5 & Cómo se complementan los métodos de la sesión matutina & los límites entre lo causal, la experiencia, la recompensa y el entorno & Expone los supuestos de investigación y los conflictos prácticos \\
6 & Cómo pasa el Context de ser ingeniería a ser aprendizaje & experience intelligence, personalización, context learning & Convierte la interacción durante el despliegue en una capacidad de largo plazo \\
7 & Por qué externalizar la capacidad & externalización de memory, tool, artifact y workflow & Reduce el costo y el riesgo de depender solo de la actualización de parámetros \\
8 & Si el Agente necesita una teoría unificada & interaction, state, tool use e intelligence boundary & Intenta definir la «segunda mitad» de la inteligencia de las máquinas \\
9 & Cómo implementarlo y gobernarlo de manera continua & retroalimentación del producto, personalización, evaluación, restricciones organizacionales y comerciales & Traduce el ciclo cerrado de investigación a sistemas reales de usuarios \\
\bottomrule
\end{longtable}

\begin{warningbox}{No equiparar automáticamente la «automejora» con ser más fuerte}
Si la retroalimentación no es verificable, la distribución del entorno es demasiado estrecha, la memoria carece de un mecanismo de olvido, o el evaluador y la política evaluada derivan juntos, el sistema puede perfectamente desviarse cada vez más a medida que aprende. Todos los diagramas de ciclo cerrado, por elegantes que sean en el curso, deben someterse a más preguntas: quién define el éxito, si la falla puede reproducirse, si la actualización puede revertirse y si se destruyen las capacidades previas.
\end{warningbox}

\IfFileExists{lecture01/lecture01-chapter.es-mx.tex}{\input{lecture01/lecture01-chapter.es-mx.tex}}{}
\IfFileExists{lecture02/lecture02-chapter.es-mx.tex}{\input{lecture02/lecture02-chapter.es-mx.tex}}{}
\IfFileExists{lecture03/lecture03-chapter.es-mx.tex}{\section{De Foresight Failure a la autoevolución con un objetivo}

El informe de Qian Cheng (钱成) estudia una capacidad de los Agentes que suele pasarse por alto: \textbf{foresight}, es decir, la capacidad de prever antes de actuar la evolución del entorno, las consecuencias de las herramientas y las restricciones de largo plazo. Un Agent puede contar con herramientas, e incluso conectarse a un world model, pero eso no garantiza que lo invoque en el momento correcto, que confíe en el modelo, o que sea capaz de convertir la predicción en mejores decisiones.

\sourcefigure{0.92}{lecture03/slides-images/slide-004-00105s.jpg}{Un Agent es un sistema interactivo capaz de percibir, razonar, actuar y recibir retroalimentación}{00:01:45--00:03:45}

\subsection{Foresight es una capacidad de un nivel superior a «saber invocar herramientas»}

En tareas complejas, el Agent necesita primero estimar los estados posteriores que puede desencadenar una acción, y solo después elegir el paso actual. El modelo del mundo, la búsqueda y los simuladores son solo herramientas que ofrecen predicciones; el verdadero foresight incluye además identificar cuándo vale la pena pagar el costo de predecir, cómo manejar la incertidumbre de la predicción, y cuándo dejar de seguir imaginando.

\sourcefigure{0.92}{lecture03/slides-images/slide-005-00225s.jpg}{El espacio de acción del Agent debe ampliarse de la invocación reactiva de herramientas a la toma de decisiones prospectiva}{00:03:45--00:05:15}

\sourcefigure{0.92}{lecture03/slides-images/slide-006-00315s.jpg}{World model como herramienta: predecir el cambio del entorno después de una acción}{00:05:15--00:07:15}

\begin{importantbox}{La brecha de capacidad no necesariamente está en la precisión del modelo}
Los experimentos muestran que, incluso si se proporciona un world model más fuerte, el Agent puede invocarlo muy pocas veces; e incluso cuando lo invoca, puede optar por ignorarlo cuando su salida entra en conflicto con el juicio previo. Por lo tanto, la falla del sistema no puede atribuirse solo a que «el predictor es impreciso»; también hay que examinar la política de invocación, la calibración de la confianza y la integración en la decisión.
\end{importantbox}

\subsection{Experimento: por qué un world model mejor no produce automáticamente un Agent más fuerte}

El informe sustituye, en varias tareas interactivas, world models con distintas capacidades, y compara el desempeño final del Agent. En el caso ideal, un predictor más fuerte debería elevar de manera monótona la tasa de éxito; sin embargo, los resultados reales muestran una ganancia limitada, lo cual indica que existe una brecha de interfaz evidente entre la capacidad del modelo y la capacidad del Agent para usar ese modelo.

\sourcefigure{0.93}{lecture03/slides-images/slide-007-00435s.jpg}{Configuración experimental: combinaciones de distintas tareas, world models y Agents}{00:07:15--00:07:45}

\sourcefigure{0.93}{lecture03/slides-images/slide-009-00540s.jpg}{El Agent presenta una reluctance significativa a invocar el world model}{00:09:00--00:09:15}

\sourcefigure{0.93}{lecture03/slides-images/slide-011-00660s.jpg}{Patrones típicos de foresight failure y pérdida de desempeño}{00:11:00--00:13:15}

Las fallas que resume el informe pueden entenderse en tres categorías. Primero, \textbf{no invocar}: el Agent actúa directamente a partir del conocimiento previo del lenguaje. Segundo, \textbf{invocar sin integrar}: la predicción se vuelve solo texto adicional, sin cambiar el plan. Tercero, \textbf{confianza errónea}: el Agent acepta la predicción sin condiciones y, en cambio, es desviado por un world model poco confiable.

\begin{warningbox}{World model no es un oracle}
Un sistema en producción debe estimar a la vez el error de predicción y el valor de uso. Si el simulador no es confiable en estados fuera de distribución, obligar a invocarlo en cada paso aumenta la latencia y propaga errores; si nunca se invoca, se desperdicia información que podría evitar fallas de alto costo. El objetivo correcto es aprender una calibrated invocation policy.
\end{warningbox}

\subsection{Foresight Governor: hacer explícito «si se piensa o no en el futuro»}

El informe propone una perspectiva de gobernanza: añadir, además del Agent de ejecución, un módulo adicional que decida cuándo solicitar una predicción, qué world model elegir, y cómo escribir el resultado de vuelta en el contexto de decisión. Este governor no debe depender solo de recordatorios temporales en el prompt, sino recibir un entrenamiento explícito, de modo que «usar foresight» se convierta en sí mismo en una conducta optimizable.

\sourcefigure{0.93}{lecture03/slides-images/slide-012-00795s.jpg}{Foresight Governor separa la invocación de la predicción de la política principal}{00:13:15--00:14:45}

\sourcefigure{0.93}{lecture03/slides-images/slide-013-00885s.jpg}{El entrenamiento explícito es una condición necesaria para obtener un foresight estable}{00:14:45--00:16:15}

\begin{knowledgebox}{Por qué no basta con un recordatorio en el prompt}
«Antes de actuar, predice con cuidado» solo modifica el contexto actual y no garantiza que el modelo mantenga la misma política ante tareas nuevas, trayectorias largas o evidencia contradictoria. El entrenamiento explícito moldea la tendencia a invocar mediante una gran cantidad de trayectorias exitosas y fallidas, y le enseña al modelo a comparar el beneficio y el costo de predecir.
\end{knowledgebox}

\subsection{Test-Time Training: convertir una sola tarea en un proceso de aprendizaje local}

El informe lleva el self-improvement un paso más allá, hasta el momento de prueba: mientras el Agent procesa la familia de tareas actual, primero genera self-awareness o experiencia reutilizable a partir de las trayectorias fallidas, y después usa esa experiencia en los episodes posteriores. La actualización puede ocurrir en la capa de memory, o bien entrar a la capa de parámetros mediante un entrenamiento a pequeña escala.

\sourcefigure{0.93}{lecture03/slides-images/slide-014-00975s.jpg}{Dos rutas del test-time training: aumento de memoria y actualización de parámetros}{00:16:15--00:18:00}

\sourcefigure{0.93}{lecture03/slides-images/slide-015-01080s.jpg}{Dos self-improvement settings: repetición de la misma tarea y transferencia de la distribución de tareas}{00:18:00--00:19:15}

\begin{importantbox}{Self-awareness y self-improvement deben distinguirse}
Self-awareness consiste en identificar «por qué fallé, qué capacidad me falta»; self-improvement consiste en cambiar, a partir de eso, la conducta posterior. Generar un texto de reflexión bien escrito no equivale a una mejora de capacidad; hay que verificar si de verdad cambió la invocación, la planeación o los resultados del entrenamiento.
\end{importantbox}

\sourcefigure{0.92}{lecture03/slides-images/slide-016-01155s.jpg}{Resultado principal: aumento de la tasa de éxito del mismo Agent tras el aprendizaje local}{00:19:15--00:20:00}

\sourcefigure{0.92}{lecture03/slides-images/slide-017-01200s.jpg}{Eficiencia de datos y crecimiento de capacidad del test-time self-improvement}{00:20:00--00:20:45}

\subsection{Generator--Solver y el curriculum automático}

Para evitar que el Agent solo repita las tareas ya existentes, el informe introduce una estructura generator--solver. El Generator genera nuevos problemas a partir del límite de capacidad del solver actual; el solver intenta resolverlos y devuelve el resultado; el diseñador de currículum ajusta después la distribución de tareas según la tasa de éxito, la dificultad y la novedad.

\sourcefigure{0.93}{lecture03/slides-images/slide-018-01245s.jpg}{Generator--Solver forma un currículum adaptativo mediante competencia y retroalimentación}{00:20:45--00:21:45}

\sourcefigure{0.93}{lecture03/slides-images/slide-020-01350s.jpg}{Curriculum reward restringe simultáneamente la resolubilidad, la dificultad y la diversidad}{00:22:30--00:23:45}

Sea $q$ la tarea generada y $p_\theta(q)$ la probabilidad de éxito del solver. Un curriculum práctico no debe limitarse a maximizar la tasa de fallas, sino preferir problemas cercanos al límite de capacidad, por ejemplo

$$
r(q)=\underbrace{p_\theta(q)(1-p_\theta(q))}_{\text{dificultad límite}}
+\lambda\,\mathrm{novelty}(q)-\mu\,\mathrm{invalid}(q).
$$

Las tareas demasiado fáciles no aportan incremento de entrenamiento, y las demasiado difíciles no ofrecen una señal aprendible; $p(1-p)$ alcanza su máximo cuando la probabilidad de éxito se acerca a $0.5$, lo cual expresa la intuición de un currículum «ligeramente por encima de la capacidad actual».

\sourcefigure{0.93}{lecture03/slides-images/slide-024-01530s.jpg}{El curriculum automático muestrea tareas cerca del límite de capacidad}{00:25:30--00:26:30}

\sourcefigure{0.93}{lecture03/slides-images/slide-026-01605s.jpg}{Resultados de mejora continua de capacidad producidos por el currículum autogenerado}{00:26:45--00:27:15}

\subsection{Resumen de la sección}

Esta sección muestra que la falla del Agent a menudo no es «carecer de world model», sino carecer de un governor que use la predicción. Por eso la autoevolución necesita entrenar explícitamente el foresight, extraer de las fallas experiencia ejecutable, y construir automáticamente, mediante generator--solver, un nuevo currículum ubicado en el límite de capacidad actual. El verdadero ciclo cerrado es: prever la falla, diagnosticar la falla, generar tareas específicas, y verificar si la capacidad cambió.
}{}
\IfFileExists{lecture04/lecture04-chapter.es-mx.tex}{\input{lecture04/lecture04-chapter.es-mx.tex}}{}
\IfFileExists{lecture05/lecture05-chapter.es-mx.tex}{\input{lecture05/lecture05-chapter.es-mx.tex}}{}
\IfFileExists{lecture06/lecture06-chapter.es-mx.tex}{\input{lecture06/lecture06-chapter.es-mx.tex}}{}
\IfFileExists{lecture07/lecture07-chapter.es-mx.tex}{\input{lecture07/lecture07-chapter.es-mx.tex}}{}
\IfFileExists{lecture08/lecture08-chapter.es-mx.tex}{\input{lecture08/lecture08-chapter.es-mx.tex}}{}
\IfFileExists{lecture09/lecture09-chapter.es-mx.tex}{\input{lecture09/lecture09-chapter.es-mx.tex}}{}
\section{Síntesis y ampliación}

La verdadera dificultad de los agentes que se automejoran no es agregar otro prompt de reflexión, sino establecer un régimen de aprendizaje capaz de operar a largo plazo: el sistema debe obtener continuamente nueva experiencia, distinguir el ruido de las regularidades transferibles, atribuir la retroalimentación a las decisiones correctas y volver a verificar las capacidades previas después de cada actualización. El modelo del mundo aporta una explicación sobre la que se puede intervenir, la síntesis de entornos aporta el currículo, la rollout infrastructure aporta la escala, la memoria externa y las herramientas aportan una persistencia de bajo costo, y la teoría y la evaluación evitan que el sistema confunda un éxito local con inteligencia general.

\subsection{Lecturas adicionales}

\begin{itemize}
    \item A Survey of Self-Evolving Agents: What, When, How, and Where to Evolve on the Path to Artificial Super Intelligence.
    \item Curious Causality-Seeking Agents in Open-ended Worlds.
    \item Agent-World: Scaling Real-World Environment Synthesis for Evolving Agents.
    \item Robust Agents Learn Causal World Models.
\end{itemize}

\end{document}
